% 中文正文，对应当前俄语版 body.tex。
\section{任务说明与确认状态}
\subsection{课题与目标}
课程设计题目：“美容沙龙搜索与服务在线预约信息系统”。

需要理发或做美甲的用户可能并不了解附近的沙龙、在那里工作的服务师及具体价格。本平台汇集多家沙龙，让用户通过地图或列表查找门店，比较服务，查看照片和评价，选择服务师与空闲时间，并创建或取消预约。

系统的目标是支持从寻找合适沙龙，到预约，再到评价实际完成的到店服务的全过程。系统面向客户、服务师、沙龙商家和平台管理员。本文中的“服务师”包括发型师、美甲师及其他美容服务专业人员。

第一阶段包括业务领域分析、需求、用例模型与描述，以及系统架构。本阶段不开发实际业务代码，也不设计数据库物理结构。

\subsection{需要与教师确认的事项}
课题、用例和技术方案目前均为提案，不代表已经获得教师批准。需要确认以下事项：
\begin{enumerate}
\item 支持多家沙龙、四类角色的平台模式，以及课程设计版本的范围；
\item 地图和列表、筛选、沙龙与服务师页面、照片、预约、取消与评价；
\item Vue、结合 Spring MVC 的 Spring Boot，以及 PostgreSQL；
\item \texttt{helios} 环境、浏览器安全访问方式、地图底图与图片存储空间。
\end{enumerate}
确认日期与结果：\textit{与教师讨论后填写}。

\section{业务领域描述}
\subsection{参与者与服务流程}
客户希望在方便的位置、预算范围内找到所需服务。做出选择需要了解地址、价格、时长、服务师、空闲时间、作品示例及评价。分散查找信息和通过电话人工确认预约，不便于比较，也可能造成排班错误。

沙龙代表，以下也称“商家”，负责注册门店，维护门店介绍、服务、价格、服务师名单、工作时段和照片。他只能管理自己的沙龙。平台为其提供公开主页和预约记录。

服务师拥有用户账号及专业资料。他可以查看自己的日程、发布作品、在必要时取消预约给自己的未来服务，并记录到店服务结果。基础版本的价格和工作时段由商家设置。

平台管理员维护公共基础数据、审核门店申请、管理访问权限，以及审核评价和图片。管理员与商家是不同角色，不负责各家沙龙的日常排班。

\subsection{地图、列表与筛选}
未设置筛选条件时，列表包含所有已在平台注册、已发布且处于启用状态的沙龙，并支持分页查看。尚未发布的申请和被封禁的门店不进入公开搜索结果。

用户可以从基础数据列表中选择地区、选择地图区域，或者使用“我附近”模式。最后一种模式会请求浏览器定位授权，并使用搜索半径，例如 1、3 或 5 千米。地理搜索模式互斥，当前模式会明确显示。拒绝定位后，用户仍可手动选择地区并使用普通列表。

通用筛选条件包括服务类型与价格范围。还可以选择日期，仅显示存在空闲时间的服务。所有条件共同生效：同一家沙龙的同一项服务必须同时满足类型与价格条件；如果选择了日期，还必须有能够提供该服务且有空闲时间的服务师。其他服务的低价不能使门店被误判为符合条件。

结果卡片显示名称、地址、照片、评分和评价数量。有距离计算起点时，还显示大致的直线距离。选择服务类型后，“起价”表示符合条件的服务中的最低价；未选择类型时，则表示该门店服务目录中的最低价，并附相应说明。没有评价时用文字标明，不显示为零分。列表可按符合条件的服务价格排序；有起点时，也可按距离排序。

切换地图和列表时保留相同筛选条件。地图标记对应所选区域内经过筛选的沙龙，不受列表当前页码限制。只有用户明确执行操作时，移动地图才会更新搜索条件。标记过多时，系统提示放大地图，并明确说明结果数量限制。

\subsection{沙龙页面与服务师页面}
沙龙页面包含详细介绍、地址和地图位置、联系方式、营业时间、附价格和时长的服务项目、服务师名单、店内环境照片、作品示例，以及门店评价和评分。选择服务后，显示适合提供该服务的服务师及可预约时间。

服务师页面包含姓名、专长、经验介绍、当前所属沙龙、可提供的服务、作品及针对该服务师的独立评价。环境照片归属于沙龙，作品照片归属于该沙龙的服务师；服务师的作品也可以在门店页面展示。

价格和时长归属于具体沙龙的服务。“美甲”等公共服务类型用于跨门店比较，并不规定统一价格。第一版中，同一家沙龙的同一项服务对所有能够提供该服务的服务师采用相同条件。

\subsection{预约、取消与评价}
客户选择沙龙、服务、服务师和日期后，系统以 15 分钟为步长提供可用的开始时间。计算时考虑完整服务时长、工作时段，以及服务师已确认的预约。整个服务必须位于同一个工作时段内；休息时间用工作时段之间的间隔表示。一项预约的结束时间可以等于下一项预约的开始时间。

界面显示空闲时间并不意味着系统为客户暂时保留该时间。确认时，服务器会再次检查可用性；针对同一服务师、时间重叠的并发请求，不得生成两条预约。价格和时长在预约中固定保存，之后修改服务目录不会改变已有预约的这些条件。

新预约的状态为“已确认”。预约开始前，客户可以取消自己的预约，服务师可以取消预约给自己的服务，商家可以取消自己沙龙的预约。服务师或商家取消时必须填写原因。取消保存后释放时间，但保留历史记录。重复取消不会改变结果，已开始的服务不得取消。

计划结束时间之后，服务师或商家将预约标记为“已完成”或“未到店”。仅仅到达结束时间，并不代表服务已经提供。在符合上述条件时，预约可从“已确认”转为“已取消”“已完成”或“未到店”；普通操作界面不允许修改终态。

对于自己已完成的预约，客户可以分别对沙龙和服务师各提交一条评价，包含 1 至 5 分的评分和可选文字。服务器通过预约确定评价对象。已取消或未到店的预约不具备评价资格。“预约—对象类型”组合的唯一性用于排除重复评价。

沙龙和服务师的评分分别根据针对各自对象、已公开的评价计算。服务师的评分不会自动计入沙龙评分。只有管理员可以附带原因隐藏评价；被隐藏的评价不计入评分，也不允许借此重复创建评价。第一版不包含评价编辑和评价回复功能。

\subsection{课程设计版本的范围}
\begin{itemize}
\item 一个城市，包含多个地区与多家沙龙；使用统一时区，价格以卢布计价。
\item 一名商家管理一家沙龙；每家沙龙有一名负责商家；每位服务师在一家沙龙工作。
\item 每个账号只有一种角色；不建立角色兼任或连锁分店模型。
\item 一条预约包含一名客户、一名服务师，以及一家沙龙的一项服务。
\item 工作日程按具体日期设置，不使用重复排班模板，不跨越午夜；可预约范围为包括当天在内的 30 天，不允许选择已经过去的时间。
\item 服务时长为正整数分钟，不必是 15 的倍数；价格非负，不设置服务师个人定价。
\item 如果修改日程、停用服务师、停用服务或解除服务师与服务的关联会影响未来预约，则拒绝该操作；应先妥善处理这些预约，例如取消。
\item 改期通过取消旧预约并创建新预约完成，不保证保留原来的时间。
\item 支持 JPEG 和 PNG 图片，每张最大 5 MB；每家沙龙最多保留 20 张启用的环境照片，每位服务师最多保留 20 张启用的作品照片。
\item 不包含在线支付、路线规划、消息推送、聊天、积分优惠、团体服务和设备管理。
\end{itemize}
必须防止的资源冲突以服务师为单位确定。系统不额外禁止客户在同一时间预约不同服务师。以上假设用于限制课程设计工作量，并需要确认。

\subsection{主要实体与关系}
\begin{longtable}{L{3.4cm}L{11.4cm}}
\toprule 实体 & 用途与主要信息 \\ \midrule
\endhead
用户 & 姓名、唯一邮箱、密码哈希、角色、访问状态。 \\
地区 & 所选城市中的地区名称；不保存复杂的地理边界。 \\
沙龙 & 商家、名称、介绍、地址、地区、坐标、联系方式、用于展示的营业时间、发布状态。 \\
服务类型 & 由管理员维护、用于比较同类服务的公共基础数据。 \\
沙龙服务 & 所属沙龙、服务类型、名称、描述、价格、时长、启用状态。 \\
服务师 & 关联用户、所属沙龙、姓名、专长、经验、启用状态。 \\
服务师—服务 & 服务师与其所属沙龙服务之间的唯一关联。 \\
工作时间 & 服务师、日期、不重叠工作时段的开始与结束时间。 \\
预约 & 客户、服务、服务师、开始和结束时间、固定保存的服务名称、价格与时长、状态，以及创建、取消和服务结果信息。 \\
照片 & 沙龙、作品所属服务师、“环境”或“作品”类型、文件路径、说明文字、排序与可见状态。 \\
评价 & 预约、对象类型、评分、文字、日期与可见状态；评价作者和具体对象通过预约确定。 \\
审核决定 & 对象、操作、管理员，以及隐藏、恢复或限制访问的时间与原因。 \\
\bottomrule
\end{longtable}
一家沙龙拥有多项服务、多位服务师和多张照片。一位服务师拥有多个工作时段和多条预约。同一家沙龙内，服务师与服务之间为多对多关系。一名客户可以拥有多条预约，每条预约最多关联两条不同对象类型的评价。

创建预约时，检查服务和服务师是否属于同一家沙龙；上传作品时，检查服务师是否属于照片关联的沙龙。页面上的营业时间用于信息展示，实际可预约时间由商家依据门店营业安排确定的工作时段决定。空闲时间选项和评分通过计算产生。已被历史记录引用的对象只归档，不进行物理删除。

\section{系统用途与自动化的预期效果}
\begin{enumerate}
\item 将地址、服务、价格、专业人员与照片集中到同一目录。
\item 支持按地点、服务、价格与空闲情况比较选择。
\item 通过作品集和已完成服务的评价，帮助客户选择服务师。
\item 自动计算可预约时间，防止同一服务师被重复预约。
\item 提供自主取消及到店历史记录。
\item 区分平台、沙龙、服务师日程和客户数据的管理权限。
\end{enumerate}
预期可以减少人工确认工作，提高选择过程的透明度。在实际部署和测量前，不宣称具体的量化经济效果。


\section{系统需求}
\subsection{功能需求}
\begin{longtable}{L{1.9cm}L{12.9cm}}
\toprule 编号 & 需求 \\ \midrule
\endhead
FR-01 & 支持客户和商家注册、登录与退出。服务师通过商家邀请获得账号；管理员不通过公开注册获得身份。 \\
FR-02 & 商家提交包含介绍、地址和坐标的沙龙申请。管理员发布该申请，或附带原因退回。 \\
FR-03 & 未登录用户可以在地图和列表上查看已注册、已发布且启用的沙龙；两种视图使用相同的筛选条件。 \\
FR-04 & 搜索支持按地区、地图区域或当前位置加半径进行。拒绝定位不会阻止手动搜索。 \\
FR-05 & 服务类型与价格条件共同生效；选择日期时，检查同一项服务的空闲时间。列表支持分页和排序。 \\
FR-06 & 沙龙页面展示介绍、地址、联系方式、营业时间、服务项目、价格、时长、服务师、环境照片、作品和沙龙评价。 \\
FR-07 & 服务师页面展示专长、经验、所属沙龙、服务项目、作品和针对服务师的评价，并提供选择时间的入口。 \\
FR-08 & 可预约时间需考虑完整服务时长、工作时段和已确认的预约，排除过去时间与重叠时段。 \\
FR-09 & 客户查看预约条件后确认预约。服务器重新检查可用性，并以原子操作防止并发冲突。 \\
FR-10 & 客户可以查看自己的预约并取消尚未开始的已确认预约；释放时间并保留历史。 \\
FR-11 & 商家只能管理自己的沙龙、服务、服务师、关联关系、工作时段和照片。拒绝会破坏未来预约有效性的修改。 \\
FR-12 & 服务师可以查看自己的日程和预约，编辑自己的介绍与作品；不得修改他人数据、价格和工作日程。 \\
FR-13 & 服务师和商家可附带原因取消属于自己管理范围的未来预约，并在计划结束时间之后标记已完成或未到店。 \\
FR-14 & 客户只能基于自己已完成的预约，分别评价沙龙和服务师；每种类型最多一条。 \\
FR-15 & 分别计算沙龙和服务师的平均评分及已公开评价数量。 \\
FR-16 & 管理员维护公共基础数据、用户访问限制与沙龙发布状态，审核评价和照片，并记录原因。 \\
FR-17 & 上传图片时，检查修改权限、实际文件格式、文件大小和数量。 \\
FR-18 & 系统说明无搜索结果、时间冲突、条件变化、权限不足和保存错误等情况，不得错误地显示操作成功。 \\
\bottomrule
\end{longtable}

\subsection{非功能需求}
下列指标是供后续验证的目标要求，并非已经完成测试所得到的结果。
\begin{longtable}{L{1.9cm}L{12.9cm}}
\toprule 编号 & 需求 \\ \midrule
\endhead
NFR-01 & 门店数据隔离：服务器检查角色、所有者及每个对象的归属。篡改标识符不得访问他人数据。 \\
NFR-02 & 一致性：并发情况下不得保存同一服务师的重叠预约；取消和状态迁移具有原子性，评价的唯一性由数据库控制。 \\
NFR-03 & 安全性：使用加盐的自适应密码哈希、服务器会话、CSRF 防护和不执行 HTML 的安全文本输出；日志中不记录机密信息。 \\
NFR-04 & 在用户主动操作并授权后请求定位。当前位置不保存在个人资料或位置历史中；提供手动选择方式。 \\
NFR-05 & 地图服务故障不影响列表与预约。定位需要浏览器支持的安全上下文。 \\
NFR-06 & 在 20 名并发用户、100 家沙龙、500 名服务师和 50\,000 条预约的规模下，搜索及可用时间查询的服务器响应时间第 95 百分位目标不超过 2 秒。外部地图和图片加载单独测量。 \\
NFR-07 & 界面使用俄语，支持宽度不小于 360 像素的屏幕，提供键盘操作和文字状态提示；切换视图保留筛选条件。 \\
NFR-08 & 仅在事务提交后返回成功确认。应用重启后预约、评价和照片仍然保留。 \\
NFR-09 & Vue、Spring Boot/Spring MVC 和 PostgreSQL 必须兼容已商定的 \texttt{helios} 演示环境。实现前明确版本与磁盘配额。 \\
NFR-10 & 在同一应用内按职责划分搜索、目录、预约、评价和数据访问；环境配置与代码分离。 \\
NFR-11 & 检查图片并使用服务器生成的文件名，不将其作为可执行文件运行；公开前移除 EXIF 元数据。图片独立于应用部署包存储。 \\
\bottomrule
\end{longtable}

\section{主要用例的模型与描述}
\subsection{参与者与系统边界}
客户使用公开搜索，登录后管理自己的预约和评价。服务师管理自己的资料和日程。商家管理自己的沙龙。管理员执行平台级操作。未登录访客表示客户尚未完成身份认证的状态，不是独立业务角色。

地图供应商属于外部技术系统。登录是受保护操作的前置条件。用例图中的关联线表示参与者参与某项操作，标有“包含”的虚线箭头表示创建预约时必须执行的检查。

\begin{figure}[htbp]
\centering
\begin{tikzpicture}[font=\small,
 uc/.style={ellipse,draw,align=center,text width=5.5cm,minimum height=0.9cm}]
\node (client) at (0,0) {客户};
\node[uc] (a) at (5.5,2.8) {用例1：注册与登录};
\node[uc] (b) at (5.5,1.4) {用例2：地图与列表搜索};
\node[uc] (c) at (5.5,0) {用例3：查看沙龙和服务师};
\node[uc] (d) at (5.5,-1.4) {用例4：选择时间并预约};
\node[uc] (e) at (5.5,-2.8) {用例5：查看和取消自己的预约};
\node[uc] (f) at (5.5,-4.2) {用例6：评价到店服务};
\node[uc] (g) at (5.5,-5.7) {用例12：检查可用性};
\foreach \u in {a,b,c,d,e,f} \draw (client.east)--(\u.west);
\draw[dashed,-{Stealth}] (d.east)--++(0.7,0)|-node[pos=0.35,right] (inclabel) {\scriptsize 包含}(g.east);
\begin{scope}[on background layer]
\node[draw,fit=(a)(g)(inclabel),inner xsep=0.35cm,inner ysep=0.45cm,label=above:{平台：客户用例}] {};
\end{scope}
\end{tikzpicture}
\caption{搜索、预约和到店评价用例模型}
\end{figure}

\begin{figure}[htbp]
\centering
\begin{tikzpicture}[font=\small,
 uc/.style={ellipse,draw,align=center,text width=5.8cm,minimum height=1cm}]
\node[align=center] (owner) at (0,2) {沙龙商家};
\node (master) at (0,-1.5) {服务师};
\node (admin) at (-0.35,-4.5) {管理员};
\node[uc] (a) at (6,3) {用例7：注册沙龙};
\node[uc] (b) at (6,1.4) {用例8：管理自己的沙龙};
\node[uc] (c) at (6,-0.2) {用例9：服务师工作};
\node[uc] (d) at (6,-1.8) {用例10：处理到店服务};
\node[uc] (e) at (6,-4.1) {用例11：平台管理与审核};
\draw (owner.east)--(a.west);
\draw (owner.east)--(b.west);
\draw (owner.east)--(d.west);
\draw (master.east)--(c.west);
\draw (master.east)--(d.west);
\draw (admin.east)--(a.west);
\draw (admin.east)--(e.west);
\begin{scope}[on background layer]
\node[draw,fit=(a)(e),inner sep=0.45cm,label=above:{平台：专业角色用例}] {};
\end{scope}
\end{tikzpicture}
\caption{沙龙商家、服务师和管理员的工作用例模型}
\end{figure}
\clearpage

\subsection{用例1：注册与登录}
\ucfield{参与者}{客户、商家；服务师和管理员仅涉及登录与退出。}
\ucfield{前置条件}{登录所用账号已存在，且未被封禁。}
\ucfield{触发事件}{打开注册或登录表单。}
\ucfield{主流程}{客户或商家输入姓名、邮箱和密码。系统检查数据并创建相应账号。登录后按角色建立会话。服务师通过商家在用例8中发出的一次性邀请设置密码；网站之外的邀请链接传递不属于本系统功能。}
\ucfield{备选流程}{邮箱已被占用、登录信息错误、邀请无效或账号被封禁时拒绝操作。退出登录时结束会话。}
\ucfield{后置条件}{账号或会话已创建；拒绝操作时权限不变。}

\subsection{用例2：地图与列表搜索}
\ucfield{参与者}{客户，无需登录。}
\ucfield{前置条件}{公开目录可用。}
\ucfield{触发事件}{打开搜索页面或修改筛选条件。}
\ucfield{主流程}{客户选择地区或区域、服务类型、价格，必要时选择日期。系统显示符合条件的卡片或地图标记。切换视图保留条件。}
\ucfield{备选流程}{执行“我附近”并授权后，以当前位置和半径搜索。拒绝授权时保留手动搜索；地图故障时使用列表。没有匹配项时提示调整条件。}
\ucfield{后置条件}{显示符合条件的结果，不创建预约。}

\subsection{用例3：查看沙龙和服务师}
\ucfield{参与者}{客户，无需登录。}
\ucfield{前置条件}{已选择已发布的沙龙或启用的服务师。}
\ucfield{触发事件}{点击卡片、地图标记或服务师姓名。}
\ucfield{主流程}{系统显示介绍、服务、价格、照片，以及针对相应对象的评价。客户选择服务和服务师，进入时间选择流程。}
\ucfield{备选流程}{明确提示没有照片或评价。已下架的沙龙不接受新预约，但已有预约仍对其相关参与者可见。}
\ucfield{后置条件}{客户获得用于选择的信息。}

\subsection{用例4：选择时间并创建预约}
\ucfield{参与者}{客户。}
\ucfield{前置条件}{已选择已发布的沙龙、启用的服务和合适的服务师；确认预约需要登录。}
\ucfield{触发事件}{选择到店日期。}
\ucfield{主流程}{系统显示空闲时间。客户选择开始时间，核对沙龙、服务、服务师、结束时间与价格。确认后，服务器在事务中执行用例12，保存预约并返回预约编号。}
\ucfield{备选流程}{已被占用的时间会被拒绝，并更新选项。价格或时长变化时需要重新确认。响应丢失时客户查看自己的预约；重复请求不会创建时间重叠的重复预约。}
\ucfield{后置条件}{成功保存一条有效预约，或者没有创建新预约。}

\subsection{用例5：查看和取消自己的预约}
\ucfield{参与者}{客户。}
\ucfield{前置条件}{客户已登录。}
\ucfield{触发事件}{打开“我的预约”。}
\ucfield{主流程}{系统显示客户自己未来和过去的预约、预约条件及状态。需要时，客户选择未来预约并确认取消。服务器检查归属和状态后保存取消结果。}
\ucfield{备选流程}{不得取消他人的预约或已经开始的预约。重复取消返回当前状态。仅查看预约不要求执行取消。}
\ucfield{后置条件}{取消时释放时间，并保留历史。}

\subsection{用例6：评价到店服务}
\ucfield{参与者}{客户。}
\ucfield{前置条件}{客户自己的预约状态为“已完成”。}
\ucfield{触发事件}{在历史记录中选择评价沙龙或服务师。}
\ucfield{主流程}{客户选择评价对象类型、评分并填写文字。服务器通过预约确定具体对象，检查唯一性并公开评价，更新对应对象的评分。}
\ucfield{备选流程}{重复的评价类型、他人的预约、已取消或未到店预约均被拒绝。同一已完成预约允许提交另一种对象类型的评价。}
\ucfield{后置条件}{保存与已完成到店服务关联的评价。}


\subsection{用例7：注册和发布沙龙}
\ucfield{参与者}{商家、管理员。}
\ucfield{前置条件}{商家已登录，且未管理其他沙龙。}
\ucfield{触发事件}{创建门店页面。}
\ucfield{主流程}{商家输入介绍、地址、地区、联系方式和营业时间，并标注坐标。系统保存申请。管理员审核并发布。}
\ucfield{备选流程}{拒绝不完整的数据。被退回的申请可以修改后重新提交。没有服务项目或工作日程的沙龙在发布后可见，但没有可预约时间。}
\ucfield{后置条件}{沙龙已发布，或仍处于不公开的申请状态。}

\subsection{用例8：管理自己的沙龙}
\ucfield{参与者}{商家。}
\ucfield{前置条件}{已登录，且沙龙属于该商家。}
\ucfield{触发事件}{打开商家管理界面。}
\ucfield{主流程}{商家编辑页面、服务、价格和照片，邀请服务师并为其分配服务项目和工作时段。服务器检查归属、数据及对未来预约的影响，然后保存修改。}
\ucfield{备选流程}{不属于本店的服务、重叠的工作日程、不符合要求的文件，以及会破坏未来预约有效性的修改均被拒绝。向已占用邮箱发送邀请会被拒绝，不会改变已有账号的角色。}
\ucfield{后置条件}{本店目录已更新；其他人的数据和已有预约中约定的条件保持不变。}

\subsection{用例9：服务师工作}
\ucfield{参与者}{服务师。}
\ucfield{前置条件}{已登录，服务师已与沙龙建立关联。}
\ucfield{触发事件}{打开个人工作界面。}
\ucfield{主流程}{服务师选择日期，查看工作时段、自己的预约、客户姓名、服务项目和时间。需要时编辑自己的介绍和作品照片。}
\ucfield{备选流程}{没有日程时显示提示。禁止修改他人的资料、价格或工作日程；不符合要求的文件被拒绝。}
\ucfield{后置条件}{显示自己的日程，并保存权限允许的修改。}

\subsection{用例10：处理到店服务}
\ucfield{参与者}{服务师、商家。}
\ucfield{前置条件}{预约属于当前服务师，或属于商家管理的沙龙。}
\ucfield{触发事件}{在预约记录中选择操作。}
\ucfield{主流程}{服务师查看自己的预约，商家查看自己沙龙的预约记录，并可按日期、服务师、客户和状态筛选。计划结束时间后，参与者将服务标记为已完成。服务器检查权限与状态，保存结果和操作者。随后客户可以评价。}
\ucfield{备选流程}{结束后可以标记未到店，此时不提供评价资格。开始前可以附带原因取消。操作时间不合适或预约已处于终态时，不允许状态迁移。}
\ucfield{后置条件}{状态与历史记录一致，评价资格与服务结果相符。}

\subsection{用例11：平台管理与审核}
\ucfield{参与者}{平台管理员。}
\ucfield{前置条件}{已使用管理员身份登录。}
\ucfield{触发事件}{选择基础数据或审核对象。}
\ucfield{主流程}{管理员维护地区和服务类型，检查内容。存在违规时，隐藏评价或照片，限制用户访问或沙龙发布状态，并保存原因和时间。}
\ucfield{备选流程}{沙龙下架后禁止新预约，但不自动取消已有预约；商家仍可处理已有预约。商家账号被封禁时，管理员可作为例外，附带原因取消受影响的未来预约。服务师被封禁后不再接受新预约，已有预约由商家处理。恢复访问或内容时，单独记录恢复决定。}
\ucfield{后置条件}{可见状态与权限已更新，被隐藏的评价不计入评分，历史保留。}

\subsection{用例12：保存前检查可用性}
\ucfield{上下文}{用例4必须包含的部分。}
\ucfield{前置条件}{事务已开始，并已保证对所选服务师相关变更进行同步控制。}
\ucfield{主流程}{服务器检查沙龙是否已发布，服务与服务师是否启用及其归属，二者的关联关系，预约条件是否最新，日期、开始时间步长、服务是否完整处于工作时段内，以及是否存在时间重叠。}
\ucfield{备选流程}{任何条件不满足都会中止保存，并给出明确原因。}
\ucfield{后置条件}{检查通过的结果仅在当前事务中有效，不能由界面的预先检查代替。}

\subsection{需求与用例的对应关系}
\begin{center}
\begin{tabular}{ll}
\toprule 需求 & 用例 \\ \midrule
FR-01 & 用例1、用例8 \\
FR-02 & 用例7 \\
FR-03--FR-05 & 用例2 \\
FR-06、FR-07 & 用例3 \\
FR-08、FR-09 & 用例4、用例12 \\
FR-10 & 用例5 \\
FR-11 & 用例8 \\
FR-12 & 用例9 \\
FR-13 & 用例10 \\
FR-14、FR-15 & 用例6、用例11 \\
FR-16 & 用例7、用例11 \\
FR-17 & 用例8、用例9 \\
FR-18 & 用例1至用例12 \\
\bottomrule
\end{tabular}
\end{center}

\section{拟采用的系统架构}
\subsection{总体方案与技术}
系统拟采用客户端—服务器架构，服务器端为单体应用。多个沙龙由同一应用和共享的 PostgreSQL 数据库提供服务，由服务器检查数据归属。在课程设计规模下，无需拆分为独立微服务，也不需要搜索集群。
\begin{itemize}
\item \textbf{Vue 3} 实现目录、筛选、沙龙与服务师页面，以及按角色划分的管理界面。
\item \textbf{Leaflet} 显示地图、标记和所选区域；地图底图由商定的供应商提供。沙龙数据来自系统自己的数据库，而非外部地图的门店目录。
\item \textbf{Spring Boot 和 Spring MVC} 提供 Java 应用和 HTTP API；Spring Security 负责会话与权限，Spring Data JPA 负责数据访问。
\item \textbf{PostgreSQL} 保存用户、沙龙、目录、工作日程、预约、评价与图片元数据。坐标存储为数值字段；在课程设计的数据规模下，距离计算不强制依赖 PostGIS。
\item \textbf{文件存储} 使用独立的持久化目录保存图片，图片路径和关联关系保存在数据库中。更新 JAR 部署包不会删除照片。
\end{itemize}
浏览器定位需要用户授权和安全上下文。地图底图供应商单独确认。使用 OpenStreetMap 公共地图瓦片时，需要遵守署名要求和服务规则，不假定它提供无限制且始终有保障的可用性。地址由商家输入，坐标通过地图选点设置，因此外部地理编码服务不是必需组件。

服务器端必须在 \texttt{helios} 上演示。Vue 静态文件可以由同一个 Spring Boot 应用提供。实现前明确 Java 和 Spring Boot 版本、浏览器安全访问方式和图片存储空间。目前不宣称已成功部署。

\begin{figure}[htbp]
\centering
\begin{tikzpicture}[font=\small,
 box/.style={draw,rounded corners,align=center,text width=10.8cm,minimum height=1.1cm},
 arrow/.style={-{Stealth},thick}]
\node[box] (ui) {浏览器：Vue、Leaflet、经授权的定位\\客户、服务师、商家、管理员};
\node[box,below=0.8cm of ui] (api) {\texttt{helios} 上的 Spring Boot / Spring MVC\\HTTP API 与 Spring Security};
\node[box,below=0.8cm of api] (logic) {目录与搜索、预约、角色工作界面、评价、照片\\检查角色及数据所属沙龙};
\node[box,below=0.8cm of logic] (data) {PostgreSQL：持久化数据与事务\\独立文件目录：照片};
\draw[arrow] (ui)--node[right]{JSON} (api);
\draw[arrow] (api)--(logic);
\draw[arrow] (logic)--(data);
\end{tikzpicture}
\caption{平台主要层次；浏览器从外部服务加载地图底图}
\end{figure}

\subsection{各组件的职责}
搜索功能将地理范围、服务、价格和日期作为一个统一请求处理。列表与地图标记由同一套服务器逻辑生成，防止结果不一致。数据量较小时，使用 PostgreSQL 筛选并根据坐标计算距离即可；外部地图负责展示。

预约服务决定时间是否可用，并在锁定服务师数据行的事务中执行检查和保存。修改工作日程和取消预约使用相同的同步机制。修改服务及沙龙发布状态也必须与创建预约协调；锁定操作采用统一顺序。详细处理协议将在实现阶段确定。仅在浏览器中检查是不够的。

评价服务检查预约归属和是否已完成，保证评价对象类型的唯一性，并分别计算评分。照片服务检查上传文件并生成安全文件名；保存失败时，不得留下指向不存在文件的公开链接。

权限服务区分客户个人数据、服务师数据和沙龙数据。请求中提供沙龙标识符，并不能证明用户具有权限。专业角色的管理界面只能获得相关到店服务所需的信息，不包含账号的保密凭据。

HTTP API 分组包括：账号与会话、搜索、沙龙与服务师页面、服务与可用时间、预约与状态、评价、照片、专业角色管理界面、审核。具体接口地址和请求结构将在下一阶段确定。

\section{结论}
新方案描述了一个支持在线预约的多沙龙搜索与比较平台。方案确定了四类角色、地图与列表统一搜索结果、筛选、门店与服务师页面、照片、基于已完成服务的评价，以及取消预约。

模型将价格归属于具体沙龙的服务，区分沙龙评价和服务师评价，并将时间冲突检查保留为预约的必要条件。系统架构采用 Vue、Spring Boot/Spring MVC 和 PostgreSQL。下一步是与教师确认项目范围、用例及演示环境。

